\begin{table}[htbp]
\centering
\footnotesize
\caption{Problema F, meia-lua (openai/gpt-4o-mini): Fase E por semente e número de exemplos.}
\label{tab:bloco2_F_meialua_faseE_nshot_gpt4mini}
\begin{tabular}{lrrrrrrrrrrrrrrr}
\toprule
 & \multicolumn{5}{c}{LLM vs real (\%)} & \multicolumn{5}{c}{LLM vs métrica (\%)} & \multicolumn{5}{c}{Perceptron vs real (\%)} \\ \cmidrule(lr){2-6} \cmidrule(lr){7-11} \cmidrule(lr){12-16}
Semente & n=0 & n=4 & n=10 & n=20 & n=40 & n=0 & n=4 & n=10 & n=20 & n=40 & n=0 & n=4 & n=10 & n=20 & n=40 \\
\midrule
42 & 55,6 & 52,6 & 60,0 & 68,1 & 65,9 & 75,6 & 62,2 & 72,6 & 86,7 & 85,9 & — & 33,3 & 69,6 & 62,2 & 58,5 \\
\midrule
\textbf{média} & \textbf{55,6} & \textbf{52,6} & \textbf{60,0} & \textbf{68,1} & \textbf{65,9} & \textbf{75,6} & \textbf{62,2} & \textbf{72,6} & \textbf{86,7} & \textbf{85,9} & \textbf{—} & \textbf{33,3} & \textbf{69,6} & \textbf{62,2} & \textbf{58,5} \\
\bottomrule
\multicolumn{16}{p{0.95\linewidth}}{\footnotesize Fase E nas 45 amostras de teste; o especialista do Problema F é o rótulo real. Perceptron baseline = perceptron linear treinado nos mesmos n exemplos. Médias dos sorteios por semente.} \\
\end{tabular}
\end{table}
